\section{For the general reader}\label{sec:general}

The prime numbers $2,3,5,7,11,\dots$ thin out as they grow, and the prime number theorem says
how fast: up to $x$ there are about $x/\log x$ of them. The Riemann hypothesis would sharpen this
estimate. It bounds the difference between the prime count and the logarithmic integral
$\operatorname{Li}(x)=\int_2^x dt/\log t$ by a constant times $\sqrt{x}\log x$ for large $x$.
Riemann reformulated the question in 1859 in terms of a function $\zeta(s)$ of a
complex variable: the error is governed by the points where $\zeta$ vanishes, and the hypothesis
says that all of its non-obvious zeros lie on one vertical line, $\operatorname{Re}s=\tfrac12$.
Its generalization to Dirichlet $L$-functions controls primes in arithmetic progressions in the
same way.

In November 2026 a short note \cite{remark} reported the author's claim to have proved both
statements and linked to Lean~4 source for the argument \cite{alpotree,lean4}. Lean is a proof
assistant: a program that checks proof terms against its rules and the axioms they use. The Lean
source contains a line that
asserts the Riemann hypothesis, in exactly the form used by Lean's standard mathematical library
Mathlib \cite{mathlib2020}, together with about 175{,}000 files of supporting mathematics.

Two questions follow, and this article keeps them apart.
\begin{enumerate}
\item \emph{Does the program accept the proof?} This is a mechanical check. An independent
  verification record reports that the theorem, in Mathlib's form, has been machine-checked by
  several independent checkers; some runs were still going (\S\ref{sec:verification}). Such a check
  relies on the checkers and on Mathlib's definitions; whether those definitions agree with the
  classical ones is a separate question, on which \S\ref{sec:verification} reports one check.
\item \emph{What is the argument?} The source is far too large to read line by line. This article
  maps the statements and proof bodies supplied in the source, explains what the statements say,
  and traces how the main pieces connect. The reading of the argument by experts remains.
\end{enumerate}

In outline the argument runs as follows. The reductions in the source lead to the assertion that no
Dirichlet $L$-function vanishes on $\operatorname{Re}s>\tfrac12$, with $s=1$ excluded for principal
characters (where these $L$-functions have a pole). To exclude such zeros, the source bounds a weighted count of primes on a sequence of
ever-larger ranges of numbers (``windows''), and uses a classical positivity principle that turns a
one-sided bound of size $\sqrt{x}$ into the absence of zeros. The bound on the prime count in turn
is reduced to bounds on certain finite groups attached to cyclotomic fields (the fields generated by
roots of unity), and the source derives those bounds from statements modelled on some of the
deepest theorems of modern number theory: the correspondence between automorphic representations and Galois representations, and the
``theta correspondence'', which transports automorphic forms between groups of different sizes. For
a von Mangoldt sum of length $X$, the note suggests a chain of about $\sqrt{X}$ theta lifts, each
mimicking the ``multiplicity'' in Stepanov's elementary proof of the Riemann hypothesis for curves
over finite fields \sttag{remark, p.~1 and item 8} \stc{bombieri-stepanov}.

Each statement in this article carries a label saying where it comes from: a published theorem
(\sttag{classical}); a location in the Lean source (\sttag{Lean}, with module replay reported in
\S\ref{sec:verification}); the author's own description in
the note (\sttag{remark}); or our reading of the source (\sttag{interpretation}).
